% =====================================================================
% DRAFT — SQ4: does the AIF variability change the perfusion maps? (AIF swap, tissue fixed)
% Methods -> sections/02_methods.tex, Results -> sections/03_results.tex,
% Discussion -> sections/04_discussion.tex. References: methods_aif_refs.bib.
%
% Show in the main text:
%   Table~\ref{tab:sq4}          tables/table_sq4.tex                 (results/combined_v1/sq4_aif_swap/table_sq4.tex)
%   Fig.~\ref{fig:sq4_change}     figures/sq4_change.png               (change per swap + within/between ratio)
%   Fig.~\ref{fig:sq4_decomp}     figures/sq4_decomposition.png        (A/B map difference = AIF part x tissue part)
% Appendix: figures/sq4_roi_examples.png, figures/sq4_ba_cross.png, figures/sq4_ba_grown.png, sq4_report.md tables.
%
% Numbers: results/combined_v1/sq4_aif_swap/sq4_report.md (QC with tail threshold 0.30, current pipeline).
% =====================================================================

% ---------------------------------------------------------------------
% METHODS
% ---------------------------------------------------------------------
\subsection{Perfusion sensitivity to the AIF (SQ4)}
\label{sec:sq4_methods}

\textbf{Rationale.} SQ1 and SQ2 showed that the AIF depends on the sampled segment and on the view. Whether that
matters depends on what it does to the perfusion maps, which are the actual output of perfDSA. In perfusion
imaging the AIF scales and shapes every parameter: a different amplitude scales CBF and CBV, and a delay or
dispersion between AIF and tissue biases CBF, MTT and Tmax~\cite{calamante2000,calamante2013}; in CT perfusion,
moving the AIF changed the maps markedly~\cite{ferreira2010}. Comparing maps of different runs mixes AIF effects with
tissue and projection effects. The AIF was therefore \emph{swapped} while the tissue curves were held fixed, so that
every difference between two maps of the same run is caused by the AIF alone.

\textbf{AIF swap.} For every run of the 68 A/B pairs the deconvolution of perfDSA was repeated on the unchanged
tissue curves (the same uniform 250\,ms sequence, which reproduces the stored maps exactly) with five AIFs: the run's
own \texttt{ica\_top} AIF (reference); its own \texttt{ica\_full\_model} and \texttt{ica\_full\_grown} AIF (segment
swaps, SQ1); the \texttt{ica\_top} AIF of the other plane of the same injection (view swap, SQ2; both planes share the
frame times, so no temporal alignment is needed); and that AIF rescaled to the peak height of the own AIF, so that only
its shape differs (shape-only view swap). A swap was analysed only when both the reference and the donor AIF were
trusted.

\textbf{Region.} No vascular-territory atlas was available. An AIF-independent parenchymal region was defined per run
from the tissue curves: pixels whose temporally smoothed enhancement exceeded about five times the frame noise and
that were not part of a large vessel or of the black collimator area (minimum-intensity value $\geq$ 120 of 170).
perfDSA sets CBV $\leq$ 0 to zero in many of these pixels; the analysis therefore used the region pixels with CBV
$>$ 0 in the reference map, and this pixel set was kept fixed for all swaps of the run, so that differences cannot
come from a changing pixel set. Per map the median CBV, CBF and MTT and the mean Tmax (more sensitive than the median
to shifts smaller than the 0.25\,s frame interval) over this region were taken. A swapped map whose median CBV or
MTT was zero (more than half of the region clipped) was counted as collapsed and left out of the log-scale
statistics.

\textbf{Analyses.} (1) Change: log ratio swapped/reference for CBV, CBF and MTT (reported as \%), difference for
Tmax; median and IQR, Wilcoxon signed-rank test, Holm correction over the four parameters of a swap. Because the view
swap acts in opposite directions on plane A and plane B tissue, all analyses were done per plane. (2) Agreement:
Bland--Altman limits of agreement~\cite{bland1999}. (3) Relevance: a change matters when it is large compared with how
much patients differ. The within-patient change (root mean square of the log ratio, which includes the bias) was
divided by the between-patient standard deviation of the log reference value in the same plane; a ratio below 1 means
that the AIF choice moves a patient less than patients differ from each other. (4) Rank stability: DSA perfusion
parameters are used as patient-level biomarkers, e.g.\ MTT and Tmax were associated with functional
outcome~\cite{vandersluijs2026}, which requires that patients keep their order. Spearman correlation between the
reference and the swapped values over the patients was computed per plane, and between the plane A and plane B maps
of the same injection. (5) Tmax: whether the earlier arrival of the grown AIF shifts Tmax was tested with the
Wilcoxon test on the Tmax shift and its Spearman correlation with the arrival difference. (6) Plane difference: the
ratio of the B and A map of one injection (each with its own AIF) was split into an AIF part (B tissue with the B-AIF
versus B tissue with the A-AIF) and a tissue part (B tissue versus A tissue with the same A-AIF).

% ---------------------------------------------------------------------
% RESULTS
% ---------------------------------------------------------------------
\subsection{Perfusion sensitivity to the AIF (SQ4)}
\label{sec:sq4_results}

\textbf{Region and swaps.} The parenchymal region covered a median 46\,\% of the image; with the reference AIF a
median of 32\,\% (IQR 21--43\,\%) of its pixels had CBV $>$ 0. Valid swaps: 44--53 per segment swap and plane, 43 per
view swap and plane. Four of about 400 swapped maps collapsed (two view swaps on A tissue, two grown swaps).

\textbf{Segment.} With the model-segmented full-ICA AIF the maps changed by 1--3\,\%
(Table~\ref{tab:sq4}, Fig.~\ref{fig:sq4_change}); the change was small compared with the patient differences (ratio
0.05--0.44) and the patients kept their order ($\rho \geq 0.85$). With the grown AIF the maps of plane B tissue
changed by $+43$\,\% (CBV), $+22$\,\% (CBF) and $+15$\,\% (MTT) ($p_\mathrm{Holm} < 0.001$); in plane A the median change
was not significant ($-8$\,\%, $-3$\,\%, $-4$\,\%) but the spread was wide (CBV IQR $-36$ to $+18$\,\%). The within/between
ratio reached 0.9--1.5 for CBV and MTT, and the patient order of CBV was largely lost ($\rho$ = 0.39 and 0.48).

\textbf{View.} Using the AIF of the other plane changed the maps most: on plane A tissue the lateral AIF lowered CBV
by 45\,\%, CBF by 37\,\% and MTT by 19\,\%; on plane B tissue the frontal AIF raised them by 73\,\%, 61\,\% and 28\,\%
(all $p_\mathrm{Holm} < 10^{-4}$). These changes exceeded the differences between patients (ratio 1.1--3.2 for CBV, CBF and
MTT) and scrambled the patient order of CBV ($\rho$ = $-0.04$ and 0.27) and MTT ($\rho$ = 0.15 and 0.44). Matching the
peak height removed most of this: the shape-only swap changed the medians by $-5$ to $+11$\,\% (CBV and MTT n.s.), the
ratio fell to 0.37--0.49 for CBF and to about 1 for CBV and MTT, and the CBF order was restored ($\rho$ = 0.94 and 0.95),
that of CBV only partly ($\rho$ = 0.69 and 0.41).

\textbf{Tmax.} Tmax was the most robust parameter: its within/between ratio stayed below 0.45 for every swap and its
patient order was preserved ($\rho \geq 0.78$). It did respond to timing: the grown AIF arrived 0.12\,s earlier
(median) and shifted the mean Tmax by $+0.19$\,s (IQR $+0.10$ to $+0.32$\,s, $p < 10^{-8}$), changing Tmax in 61\,\% of the
pixels; the shift correlated with the arrival difference ($\rho = 0.52$). The view swaps shifted Tmax by about
$\pm$0.13\,s.

\textbf{Plane A versus plane B maps.} For the same injection the lateral map had a lower CBV (B/A 0.38), CBF (0.66)
and MTT (0.57) than the frontal map (over the fixed parenchymal region; over the perfused pixels of each map, as in
SQ2, CBV B/A was 0.42), and the patient order agreed only for CBF and Tmax ($\rho$ = 0.79 and 0.80), not for
CBV and MTT ($\rho$ = 0.35 and 0.35). Split into its parts (Fig.~\ref{fig:sq4_decomp}), the CBF difference came
entirely from the AIF (AIF part 0.62, tissue part 1.05), whereas the CBV and MTT differences came from both the AIF
(0.58 and 0.78) and the tissue (0.77 and 0.73).

\begin{table*}[t]
\centering
\caption{Effect of the AIF on the perfusion maps with the tissue curves held fixed, for plane A and plane B tissue. Change: median (IQR) of the swapped map relative to the map with the run's own \texttt{ica\_top} AIF over a fixed parenchymal region (Tmax: shift of the region mean, s); $^\ast p_\mathrm{Holm} < 0.05$ (Wilcoxon, four parameters per swap and plane). Ratio: within-patient change (RMS) / between-patient SD. $\rho$: Spearman rank correlation of the patients between reference and swap.}
\label{tab:sq4}
\footnotesize
\begin{tabular}{ll rrr rrr}
\hline
 & & \multicolumn{3}{c}{Plane A tissue} & \multicolumn{3}{c}{Plane B tissue} \\
AIF swap & Param. & Change (IQR) & Ratio & $\rho$ & Change (IQR) & Ratio & $\rho$ \\
\hline
segment: ica_full_model & CBV & -2\% (-6\%, +0\%)$^\ast$ & 0.17 & 0.97 & -3\% (-4\%, +0\%)$^\ast$ & 0.44 & 0.85 \\
 & CBF & -2\% (-4\%, -0\%)$^\ast$ & 0.12 & 0.99 & -2\% (-3\%, -1\%)$^\ast$ & 0.23 & 0.97 \\
 & MTT & -1\% (-4\%, +0\%)$^\ast$ & 0.19 & 0.97 & -1\% (-3\%, +1\%)$^\ast$ & 0.28 & 0.98 \\
 & Tmax & +0.00 (-0.01, +0.01) & 0.08 & 0.99 & +0.01 (+0.00, +0.01)$^\ast$ & 0.05 & 1.00 \\
\hline
segment: ica_full_grown & CBV & -8\% (-36\%, +18\%) & 0.91 & 0.39 & +42\% (+17\%, +67\%)$^\ast$ & 1.54 & 0.48 \\
 & CBF & -3\% (-17\%, +7\%) & 0.46 & 0.73 & +22\% (+10\%, +48\%)$^\ast$ & 0.73 & 0.91 \\
 & MTT & -4\% (-21\%, +9\%) & 1.23 & 0.45 & +15\% (+1\%, +27\%)$^\ast$ & 0.92 & 0.70 \\
 & Tmax & +0.24 (+0.10, +0.36)$^\ast$ & 0.43 & 0.95 & +0.17 (+0.09, +0.24)$^\ast$ & 0.32 & 0.97 \\
\hline
view: other plane's ica_top & CBV & -45\% (-70\%, -26\%)$^\ast$ & 1.55 & -0.04 & +73\% (+36\%, +205\%)$^\ast$ & 3.17 & 0.27 \\
 & CBF & -37\% (-44\%, -19\%)$^\ast$ & 1.13 & 0.59 & +61\% (+33\%, +82\%)$^\ast$ & 1.34 & 0.77 \\
 & MTT & -19\% (-44\%, -3\%)$^\ast$ & 1.53 & 0.15 & +28\% (+1\%, +78\%)$^\ast$ & 1.85 & 0.44 \\
 & Tmax & -0.14 (-0.21, +0.03)$^\ast$ & 0.28 & 0.79 & +0.12 (-0.04, +0.19)$^\ast$ & 0.34 & 0.93 \\
\hline
view, shape only (peak-matched) & CBV & -2\% (-23\%, +10\%) & 0.97 & 0.69 & +10\% (-15\%, +36\%) & 1.01 & 0.41 \\
 & CBF & -4\% (-13\%, +2\%)$^\ast$ & 0.49 & 0.94 & +7\% (-1\%, +19\%)$^\ast$ & 0.37 & 0.95 \\
 & MTT & -4\% (-16\%, +13\%) & 1.12 & 0.56 & +5\% (-14\%, +18\%) & 0.88 & 0.72 \\
 & Tmax & -0.15 (-0.21, +0.02)$^\ast$ & 0.40 & 0.93 & +0.13 (+0.04, +0.18)$^\ast$ & 0.27 & 0.95 \\
\hline
\end{tabular}
\end{table*}


\begin{figure*}[t]
\centering
\includegraphics[width=\textwidth]{figures/sq4_change.png}
\caption{Effect of the AIF on the perfusion maps with the tissue curves held fixed, for plane A (left) and plane B
(right) tissue. Top: change relative to the map with the run's own \texttt{ica\_top} AIF (median, IQR; CBV, CBF, MTT in
\%, Tmax in 0.01\,s). Bottom: within-patient change divided by the between-patient spread; above the dashed line the AIF
choice moves a patient more than patients differ from each other.}
\label{fig:sq4_change}
\end{figure*}

\begin{figure}[t]
\centering
\includegraphics[width=\columnwidth]{figures/sq4_decomposition.png}
\caption{Ratio of the plane B and plane A perfusion map of the same injection (purple), split into an AIF part (B tissue
with the B-AIF vs.\ the A-AIF, orange) and a tissue part (B vs.\ A tissue with the same A-AIF, blue); median and IQR,
43 injections.}
\label{fig:sq4_decomp}
\end{figure}

% ---------------------------------------------------------------------
% DISCUSSION
% ---------------------------------------------------------------------
\subsection{Does the AIF variability change the maps?}
\label{sec:sq4_discussion}

\textbf{Yes, for realistic alternatives, and as much as patients differ.} Holding the tissue fixed isolates the AIF:
any difference between two maps of the same run is caused by the AIF choice. A near-identical AIF (model-segmented
full ICA) changes the maps by a few per cent and leaves the patient order intact, so small segmentation differences
are harmless. A different segment (grown full ICA) or the AIF of the other view changes CBV, CBF and MTT by tens of
per cent, by as much as or more than patients differ from each other, and reorders the patients for CBV and MTT. For a
biomarker that is meant to separate patients, such as MTT and Tmax for functional outcome~\cite{vandersluijs2026}, the
AIF choice is therefore not a detail: it can change which patients appear abnormal.

\textbf{How: mostly amplitude, partly shape and timing.} The deconvolution makes CBF and CBV inversely proportional to
the AIF amplitude, and the view swap mainly changes the amplitude (SQ2: the lateral AIF is about a third higher because
of the exposure). Matching the peak height removed the systematic part of the view effect and restored the CBF order,
which confirms that amplitude is the main path. What remains after matching, a per-patient spread in CBV and MTT
comparable to the patient differences, is due to the shape of the curve: CBV depends on the AIF area and MTT on its
width, and these differ between views even at equal peak height. Timing acts on Tmax: the grown AIF, sampled upstream,
arrives about 0.12\,s earlier and shifts Tmax by about 0.19\,s, the delay effect described for MR and CT
perfusion~\cite{calamante2000}. Because the frame interval is 0.25\,s, this shift is close to one frame and changes the
Tmax of most pixels by one step, yet it is small compared with the patient differences, so Tmax ranks are preserved.

\textbf{What this adds to SQ2.} The CBF of the lateral and frontal map of one injection differs entirely because of
the AIF: with the same AIF the two views give the same CBF (tissue part 1.05). For CBV and MTT a tissue part remains:
with the same AIF the lateral tissue still yields about a quarter lower CBV and MTT. The view difference in CBV and MTT is
therefore not only an AIF problem; it also reflects how the parenchyma is projected, which is consistent with the lower
tissue signal found in SQ2.

\textbf{Limitations.} (i) The region is the whole enhancing parenchyma of a view, not a vascular territory, and
includes contralateral parenchyma in runs with cross-filling. (ii) perfDSA sets CBV $\leq$ 0 to zero in about two thirds
of the enhancing pixels; the medians describe the remaining third, which is fixed per run but is itself defined by the
reference AIF. (iii) Tmax is quantised at 0.25\,s; the region mean is used to detect smaller shifts. (iv) There is no
reference standard: the swaps show how much the maps depend on the AIF, not which AIF gives the correct map. (v)
Between-patient spread is estimated from 43--54 runs per plane and includes residual measurement noise, so the ratio is
a conservative comparison. (vi) QC used the tail threshold of 0.30. (vii) The decomposition uses medians over regions
that differ between the views, so the parts are approximate (their product is not exactly the total).

\textbf{Recommendation.} The AIF should be defined in one fixed way, sampled from the supraclinoid ICA as perfDSA
does, and taken from the same view whenever patients or time points are compared, because maps made with another
segment or view are not comparable. Small differences in the outline of the supraclinoid mask are harmless, so quality
control is better spent on excluding wrong masks and on checking vessel overlap and exposure than on refining the
segmentation. When perfusion has to be compared across acquisitions, Tmax and CBF are the most reliable parameters:
Tmax does not depend on the AIF amplitude, and CBF agrees between views once the AIF amplitude is matched, whereas CBV and
MTT remain sensitive to the shape of the AIF and to the projection. Where absolute values cannot be standardised, relative
parameters within one map, normalised to a reference region as is common in CT and MR
perfusion~\cite{copen2014,delarosa2023}, remove a common AIF factor altogether.

\textbf{Conclusion.} The variability of the AIF changes the perfusion maps. A practically identical AIF changes them by a
few per cent, but a different ICA segment or the AIF of the other view changes CBV, CBF and MTT by tens of per cent, as
much as or more than patients differ from each other, and changes the order of the patients for CBV and MTT. The amplitude
of the AIF is the main cause; its shape adds patient-specific error in CBV and MTT, and its timing shifts Tmax. Tmax is the
most robust parameter. The AIF choice is therefore part of the measurement and must be fixed and reported for DSA
perfusion values to be interpretable.
